% Methodology description for scientific publication
% Assurance Gaps in an Integrated Safety and Cybersecurity Case
% for an AI-Based Perception Component in Highly Automated Driving
%
% This file provides the detailed methodology description that can be
% included directly in a LaTeX paper using % Methodology description for scientific publication
% Assurance Gaps in an Integrated Safety and Cybersecurity Case
% for an AI-Based Perception Component in Highly Automated Driving
%
% This file provides the detailed methodology description that can be
% included directly in a LaTeX paper using % Methodology description for scientific publication
% Assurance Gaps in an Integrated Safety and Cybersecurity Case
% for an AI-Based Perception Component in Highly Automated Driving
%
% This file provides the detailed methodology description that can be
% included directly in a LaTeX paper using % Methodology description for scientific publication
% Assurance Gaps in an Integrated Safety and Cybersecurity Case
% for an AI-Based Perception Component in Highly Automated Driving
%
% This file provides the detailed methodology description that can be
% included directly in a LaTeX paper using \input{docs/latex/methodology.tex}
% Requires: booktabs, amsmath, hyperref packages

\section{Methodology}\label{sec:methodology}

The analysis follows a five-step constructive integration methodology.
The input is a set of applicable standards; the output is an integrated
GSN argument pattern with identified inconsistencies and gaps.
Each step produces a verifiable intermediate artefact.

\subsection{System Definition and Standards Selection}\label{sec:system-def}

The case study component is a LiDAR-based 3D object detection module for
an SAE Level~4+ highly automated driving system. The architecture follows
the SECOND detector~\cite{Yan2018} with a deep ensemble of five
independently trained instances for uncertainty
quantification~\cite{Lakshminarayanan2017}.

The applicable standards are selected based on the following criteria:
(i)~the standard addresses at least one aspect of the component's safety
or cybersecurity assurance, (ii)~the standard is published (not draft),
and (iii)~the standard contains normative requirements or provides the
base assurance argument pattern. Five standards satisfy these criteria
(Table~\ref{tab:standards}). ISO/IEC~TR~5469 is included as informative
supplementary guidance.

The system receives an ASIL~D assignment from the hazard analysis and
risk assessment (HARA) per ISO~26262-3 Clause~6: severity S3 (life-threatening
injuries), exposure E4 (high probability in urban driving), controllability C3
(no human driver at SAE Level~4+). The threat analysis and risk assessment
(TARA) per ISO/SAE~21434 Clause~15 identifies three threat scenarios:
LiDAR spoofing (injected ghost points), adversarial point cloud perturbation,
and model poisoning through compromised training data.

\subsection{Step 1: Claim Extraction}\label{sec:step1}

For each applicable standard, we extract \emph{claims}: normative
requirements (\texttt{shall} statements), defined work products, and
explicitly required lifecycle activities that an assurance argument for
the AI-based perception component must address.

Each claim is assigned:
\begin{itemize}
    \item a unique identifier (e.g., \texttt{CLM-26262-HARA-01}),
    \item the source clause reference,
    \item a claim type (normative requirement, work product, or lifecycle activity),
    \item a target GSN goal node (determined in Step~3),
    \item a primary lifecycle phase (determined in Step~2).
\end{itemize}

Informative notes and recommendations are recorded but not treated as
normative claims. ISO/IEC~TR~5469, being a Technical Report, contributes
informative guidance only.

\subsection{Step 2: Lifecycle-Phase Mapping}\label{sec:step2}

Claims and clauses are mapped to six lifecycle phases that cover the
complete lifecycle of an AI-based perception component:
\begin{enumerate}
    \item \textbf{Concept / Requirements}: item definition, HARA, TARA,
          functional safety concept, SOTIF specification.
    \item \textbf{Design / Training}: AI model architecture, training process,
          cybersecurity control design, SOTIF design measures.
    \item \textbf{Verification \& Validation}: MC/DC coverage, scenario-based
          testing, AI-specific evaluation, penetration testing.
    \item \textbf{Integration / Deployment}: system integration, pre-deployment
          validation, release decision.
    \item \textbf{Operation / Monitoring}: runtime OOD detection, SOTIF field
          monitoring, cybersecurity incident monitoring.
    \item \textbf{Modification / Re-assurance}: OTA update management,
          change impact analysis, re-validation.
\end{enumerate}

The mapping produces a coverage matrix (Table~\ref{tab:coverage}) showing
how many clauses from each standard apply to each phase. Phases with low
coverage from any single standard are candidates for gaps.

\subsection{Step 3: GSN Construction}\label{sec:step3}

The integrated GSN argument pattern is constructed by:
\begin{enumerate}
    \item Starting with ISO/PAS~8800 Annex~B as the base structure. This
          provides goals G1--G6, strategy S1, and assumption A1.4.
    \item Adding context nodes: C1 (ASIL~D from HARA, ISO~26262-3 Clause~6)
          and C2 (TARA results, ISO/SAE~21434 Clause~15).
    \item Augmenting existing goals G2, G4, G5, G6 with claims from
          ISO~21448 and ISO/SAE~21434. The base pattern's G5 (V\&V)
          receives claims from all four normative standards.
    \item Adding new goals G7 (SOTIF residual risk acceptance, from
          ISO~21448 Clause~6.5) and G8 (cybersecurity risk management,
          from ISO/SAE~21434 Clause~3.1.11).
    \item Identifying G9 (modification assurance) as an undeveloped goal
          where no standard prescribes a complete re-assurance workflow
          for AI model modifications.
\end{enumerate}

The resulting structure (Table~\ref{tab:gsn-goals}) extends the base
pattern from 6 to 9 goals. Strategy S1 is reformulated to argue over
four assurance domains rather than one.

\subsection{Step 4: Junction-Point Analysis}\label{sec:step4}

A \emph{junction point} is a goal node where claims from two or more
standards meet. At each junction point, we analyse whether the standards
impose compatible requirements. Three types of inconsistency are
distinguished:
\begin{itemize}
    \item \textbf{Structural (S)}: different risk classification frameworks,
          evidence structures, or architectural assumptions.
    \item \textbf{Terminological (T)}: different terms for the same concept
          or the same term with different definitions.
    \item \textbf{Methodological (M)}: different methods prescribed for the
          same assurance objective.
\end{itemize}

The analysis identifies seven inconsistencies (Table~\ref{tab:inconsistencies}).
The central finding is I-2 (evidence type asymmetry at G5): four fundamentally
different evidence types---MC/DC structural coverage (ISO~26262), scenario-based
testing (ISO~21448), AI uncertainty quantification (ISO/PAS~8800), and
penetration testing (ISO/SAE~21434)---must support a single sufficiency
claim, yet no standard defines how they combine.

\subsection{Step 5: Gap Identification}\label{sec:step5}

Gaps are identified where the integrated argument requires a claim or
evidence that no applicable standard provides. Three gap types are
distinguished:
\begin{itemize}
    \item \textbf{Missing Claim (MC)}: the argument requires a claim
          but no standard contains it.
    \item \textbf{Missing Evidence (ME)}: a claim exists but no standard
          prescribes how to produce evidence for AI components.
    \item \textbf{Unresolved Inconsistency (UI)}: an inconsistency from
          Step~4 requires resolution that no standard provides.
\end{itemize}

Gaps are further classified as \emph{integration-induced} if they appear
only when standards are combined (not visible when each standard is applied
independently). The analysis identifies six gaps (Table~\ref{tab:gaps}),
of which two (Gap-3 and Gap-4) are integration-induced.

\subsection{Evaluation Protocol}\label{sec:evaluation}

The case study evaluation uses the CARLA simulator~\cite{Dosovitskiy2017}
under parametrically controlled weather conditions. Rain intensity varies
across five levels (0, 10, 25, 50, 100~mm/h) and fog visibility across
five levels (500, 200, 100, 50, 10~m), producing a grid of 25 weather
combinations. Each combination is evaluated over 50 scenes with randomised
traffic configurations (seed: 42).

A weather condition is classified as a \emph{SOTIF triggering condition}
(per ISO~21448 Clause~7) if rain exceeds 20~mm/h (reducing LiDAR point
density through absorption and scattering) or visibility falls below 200~m
(causing spurious reflections and reduced sensor range).

For each scene, we record: number of ground-truth objects, detections
(true positives, false positives, false negatives), mean geometric
divergence across ensemble members, and the number of uncertain detections
(divergence exceeding 0.5). These metrics map to the four evidence types
at G5 (Table~\ref{tab:evidence-types}).

% Requires: booktabs, amsmath, hyperref packages

\section{Methodology}\label{sec:methodology}

The analysis follows a five-step constructive integration methodology.
The input is a set of applicable standards; the output is an integrated
GSN argument pattern with identified inconsistencies and gaps.
Each step produces a verifiable intermediate artefact.

\subsection{System Definition and Standards Selection}\label{sec:system-def}

The case study component is a LiDAR-based 3D object detection module for
an SAE Level~4+ highly automated driving system. The architecture follows
the SECOND detector~\cite{Yan2018} with a deep ensemble of five
independently trained instances for uncertainty
quantification~\cite{Lakshminarayanan2017}.

The applicable standards are selected based on the following criteria:
(i)~the standard addresses at least one aspect of the component's safety
or cybersecurity assurance, (ii)~the standard is published (not draft),
and (iii)~the standard contains normative requirements or provides the
base assurance argument pattern. Five standards satisfy these criteria
(Table~\ref{tab:standards}). ISO/IEC~TR~5469 is included as informative
supplementary guidance.

The system receives an ASIL~D assignment from the hazard analysis and
risk assessment (HARA) per ISO~26262-3 Clause~6: severity S3 (life-threatening
injuries), exposure E4 (high probability in urban driving), controllability C3
(no human driver at SAE Level~4+). The threat analysis and risk assessment
(TARA) per ISO/SAE~21434 Clause~15 identifies three threat scenarios:
LiDAR spoofing (injected ghost points), adversarial point cloud perturbation,
and model poisoning through compromised training data.

\subsection{Step 1: Claim Extraction}\label{sec:step1}

For each applicable standard, we extract \emph{claims}: normative
requirements (\texttt{shall} statements), defined work products, and
explicitly required lifecycle activities that an assurance argument for
the AI-based perception component must address.

Each claim is assigned:
\begin{itemize}
    \item a unique identifier (e.g., \texttt{CLM-26262-HARA-01}),
    \item the source clause reference,
    \item a claim type (normative requirement, work product, or lifecycle activity),
    \item a target GSN goal node (determined in Step~3),
    \item a primary lifecycle phase (determined in Step~2).
\end{itemize}

Informative notes and recommendations are recorded but not treated as
normative claims. ISO/IEC~TR~5469, being a Technical Report, contributes
informative guidance only.

\subsection{Step 2: Lifecycle-Phase Mapping}\label{sec:step2}

Claims and clauses are mapped to six lifecycle phases that cover the
complete lifecycle of an AI-based perception component:
\begin{enumerate}
    \item \textbf{Concept / Requirements}: item definition, HARA, TARA,
          functional safety concept, SOTIF specification.
    \item \textbf{Design / Training}: AI model architecture, training process,
          cybersecurity control design, SOTIF design measures.
    \item \textbf{Verification \& Validation}: MC/DC coverage, scenario-based
          testing, AI-specific evaluation, penetration testing.
    \item \textbf{Integration / Deployment}: system integration, pre-deployment
          validation, release decision.
    \item \textbf{Operation / Monitoring}: runtime OOD detection, SOTIF field
          monitoring, cybersecurity incident monitoring.
    \item \textbf{Modification / Re-assurance}: OTA update management,
          change impact analysis, re-validation.
\end{enumerate}

The mapping produces a coverage matrix (Table~\ref{tab:coverage}) showing
how many clauses from each standard apply to each phase. Phases with low
coverage from any single standard are candidates for gaps.

\subsection{Step 3: GSN Construction}\label{sec:step3}

The integrated GSN argument pattern is constructed by:
\begin{enumerate}
    \item Starting with ISO/PAS~8800 Annex~B as the base structure. This
          provides goals G1--G6, strategy S1, and assumption A1.4.
    \item Adding context nodes: C1 (ASIL~D from HARA, ISO~26262-3 Clause~6)
          and C2 (TARA results, ISO/SAE~21434 Clause~15).
    \item Augmenting existing goals G2, G4, G5, G6 with claims from
          ISO~21448 and ISO/SAE~21434. The base pattern's G5 (V\&V)
          receives claims from all four normative standards.
    \item Adding new goals G7 (SOTIF residual risk acceptance, from
          ISO~21448 Clause~6.5) and G8 (cybersecurity risk management,
          from ISO/SAE~21434 Clause~3.1.11).
    \item Identifying G9 (modification assurance) as an undeveloped goal
          where no standard prescribes a complete re-assurance workflow
          for AI model modifications.
\end{enumerate}

The resulting structure (Table~\ref{tab:gsn-goals}) extends the base
pattern from 6 to 9 goals. Strategy S1 is reformulated to argue over
four assurance domains rather than one.

\subsection{Step 4: Junction-Point Analysis}\label{sec:step4}

A \emph{junction point} is a goal node where claims from two or more
standards meet. At each junction point, we analyse whether the standards
impose compatible requirements. Three types of inconsistency are
distinguished:
\begin{itemize}
    \item \textbf{Structural (S)}: different risk classification frameworks,
          evidence structures, or architectural assumptions.
    \item \textbf{Terminological (T)}: different terms for the same concept
          or the same term with different definitions.
    \item \textbf{Methodological (M)}: different methods prescribed for the
          same assurance objective.
\end{itemize}

The analysis identifies seven inconsistencies (Table~\ref{tab:inconsistencies}).
The central finding is I-2 (evidence type asymmetry at G5): four fundamentally
different evidence types---MC/DC structural coverage (ISO~26262), scenario-based
testing (ISO~21448), AI uncertainty quantification (ISO/PAS~8800), and
penetration testing (ISO/SAE~21434)---must support a single sufficiency
claim, yet no standard defines how they combine.

\subsection{Step 5: Gap Identification}\label{sec:step5}

Gaps are identified where the integrated argument requires a claim or
evidence that no applicable standard provides. Three gap types are
distinguished:
\begin{itemize}
    \item \textbf{Missing Claim (MC)}: the argument requires a claim
          but no standard contains it.
    \item \textbf{Missing Evidence (ME)}: a claim exists but no standard
          prescribes how to produce evidence for AI components.
    \item \textbf{Unresolved Inconsistency (UI)}: an inconsistency from
          Step~4 requires resolution that no standard provides.
\end{itemize}

Gaps are further classified as \emph{integration-induced} if they appear
only when standards are combined (not visible when each standard is applied
independently). The analysis identifies six gaps (Table~\ref{tab:gaps}),
of which two (Gap-3 and Gap-4) are integration-induced.

\subsection{Evaluation Protocol}\label{sec:evaluation}

The case study evaluation uses the CARLA simulator~\cite{Dosovitskiy2017}
under parametrically controlled weather conditions. Rain intensity varies
across five levels (0, 10, 25, 50, 100~mm/h) and fog visibility across
five levels (500, 200, 100, 50, 10~m), producing a grid of 25 weather
combinations. Each combination is evaluated over 50 scenes with randomised
traffic configurations (seed: 42).

A weather condition is classified as a \emph{SOTIF triggering condition}
(per ISO~21448 Clause~7) if rain exceeds 20~mm/h (reducing LiDAR point
density through absorption and scattering) or visibility falls below 200~m
(causing spurious reflections and reduced sensor range).

For each scene, we record: number of ground-truth objects, detections
(true positives, false positives, false negatives), mean geometric
divergence across ensemble members, and the number of uncertain detections
(divergence exceeding 0.5). These metrics map to the four evidence types
at G5 (Table~\ref{tab:evidence-types}).

% Requires: booktabs, amsmath, hyperref packages

\section{Methodology}\label{sec:methodology}

The analysis follows a five-step constructive integration methodology.
The input is a set of applicable standards; the output is an integrated
GSN argument pattern with identified inconsistencies and gaps.
Each step produces a verifiable intermediate artefact.

\subsection{System Definition and Standards Selection}\label{sec:system-def}

The case study component is a LiDAR-based 3D object detection module for
an SAE Level~4+ highly automated driving system. The architecture follows
the SECOND detector~\cite{Yan2018} with a deep ensemble of five
independently trained instances for uncertainty
quantification~\cite{Lakshminarayanan2017}.

The applicable standards are selected based on the following criteria:
(i)~the standard addresses at least one aspect of the component's safety
or cybersecurity assurance, (ii)~the standard is published (not draft),
and (iii)~the standard contains normative requirements or provides the
base assurance argument pattern. Five standards satisfy these criteria
(Table~\ref{tab:standards}). ISO/IEC~TR~5469 is included as informative
supplementary guidance.

The system receives an ASIL~D assignment from the hazard analysis and
risk assessment (HARA) per ISO~26262-3 Clause~6: severity S3 (life-threatening
injuries), exposure E4 (high probability in urban driving), controllability C3
(no human driver at SAE Level~4+). The threat analysis and risk assessment
(TARA) per ISO/SAE~21434 Clause~15 identifies three threat scenarios:
LiDAR spoofing (injected ghost points), adversarial point cloud perturbation,
and model poisoning through compromised training data.

\subsection{Step 1: Claim Extraction}\label{sec:step1}

For each applicable standard, we extract \emph{claims}: normative
requirements (\texttt{shall} statements), defined work products, and
explicitly required lifecycle activities that an assurance argument for
the AI-based perception component must address.

Each claim is assigned:
\begin{itemize}
    \item a unique identifier (e.g., \texttt{CLM-26262-HARA-01}),
    \item the source clause reference,
    \item a claim type (normative requirement, work product, or lifecycle activity),
    \item a target GSN goal node (determined in Step~3),
    \item a primary lifecycle phase (determined in Step~2).
\end{itemize}

Informative notes and recommendations are recorded but not treated as
normative claims. ISO/IEC~TR~5469, being a Technical Report, contributes
informative guidance only.

\subsection{Step 2: Lifecycle-Phase Mapping}\label{sec:step2}

Claims and clauses are mapped to six lifecycle phases that cover the
complete lifecycle of an AI-based perception component:
\begin{enumerate}
    \item \textbf{Concept / Requirements}: item definition, HARA, TARA,
          functional safety concept, SOTIF specification.
    \item \textbf{Design / Training}: AI model architecture, training process,
          cybersecurity control design, SOTIF design measures.
    \item \textbf{Verification \& Validation}: MC/DC coverage, scenario-based
          testing, AI-specific evaluation, penetration testing.
    \item \textbf{Integration / Deployment}: system integration, pre-deployment
          validation, release decision.
    \item \textbf{Operation / Monitoring}: runtime OOD detection, SOTIF field
          monitoring, cybersecurity incident monitoring.
    \item \textbf{Modification / Re-assurance}: OTA update management,
          change impact analysis, re-validation.
\end{enumerate}

The mapping produces a coverage matrix (Table~\ref{tab:coverage}) showing
how many clauses from each standard apply to each phase. Phases with low
coverage from any single standard are candidates for gaps.

\subsection{Step 3: GSN Construction}\label{sec:step3}

The integrated GSN argument pattern is constructed by:
\begin{enumerate}
    \item Starting with ISO/PAS~8800 Annex~B as the base structure. This
          provides goals G1--G6, strategy S1, and assumption A1.4.
    \item Adding context nodes: C1 (ASIL~D from HARA, ISO~26262-3 Clause~6)
          and C2 (TARA results, ISO/SAE~21434 Clause~15).
    \item Augmenting existing goals G2, G4, G5, G6 with claims from
          ISO~21448 and ISO/SAE~21434. The base pattern's G5 (V\&V)
          receives claims from all four normative standards.
    \item Adding new goals G7 (SOTIF residual risk acceptance, from
          ISO~21448 Clause~6.5) and G8 (cybersecurity risk management,
          from ISO/SAE~21434 Clause~3.1.11).
    \item Identifying G9 (modification assurance) as an undeveloped goal
          where no standard prescribes a complete re-assurance workflow
          for AI model modifications.
\end{enumerate}

The resulting structure (Table~\ref{tab:gsn-goals}) extends the base
pattern from 6 to 9 goals. Strategy S1 is reformulated to argue over
four assurance domains rather than one.

\subsection{Step 4: Junction-Point Analysis}\label{sec:step4}

A \emph{junction point} is a goal node where claims from two or more
standards meet. At each junction point, we analyse whether the standards
impose compatible requirements. Three types of inconsistency are
distinguished:
\begin{itemize}
    \item \textbf{Structural (S)}: different risk classification frameworks,
          evidence structures, or architectural assumptions.
    \item \textbf{Terminological (T)}: different terms for the same concept
          or the same term with different definitions.
    \item \textbf{Methodological (M)}: different methods prescribed for the
          same assurance objective.
\end{itemize}

The analysis identifies seven inconsistencies (Table~\ref{tab:inconsistencies}).
The central finding is I-2 (evidence type asymmetry at G5): four fundamentally
different evidence types---MC/DC structural coverage (ISO~26262), scenario-based
testing (ISO~21448), AI uncertainty quantification (ISO/PAS~8800), and
penetration testing (ISO/SAE~21434)---must support a single sufficiency
claim, yet no standard defines how they combine.

\subsection{Step 5: Gap Identification}\label{sec:step5}

Gaps are identified where the integrated argument requires a claim or
evidence that no applicable standard provides. Three gap types are
distinguished:
\begin{itemize}
    \item \textbf{Missing Claim (MC)}: the argument requires a claim
          but no standard contains it.
    \item \textbf{Missing Evidence (ME)}: a claim exists but no standard
          prescribes how to produce evidence for AI components.
    \item \textbf{Unresolved Inconsistency (UI)}: an inconsistency from
          Step~4 requires resolution that no standard provides.
\end{itemize}

Gaps are further classified as \emph{integration-induced} if they appear
only when standards are combined (not visible when each standard is applied
independently). The analysis identifies six gaps (Table~\ref{tab:gaps}),
of which two (Gap-3 and Gap-4) are integration-induced.

\subsection{Evaluation Protocol}\label{sec:evaluation}

The case study evaluation uses the CARLA simulator~\cite{Dosovitskiy2017}
under parametrically controlled weather conditions. Rain intensity varies
across five levels (0, 10, 25, 50, 100~mm/h) and fog visibility across
five levels (500, 200, 100, 50, 10~m), producing a grid of 25 weather
combinations. Each combination is evaluated over 50 scenes with randomised
traffic configurations (seed: 42).

A weather condition is classified as a \emph{SOTIF triggering condition}
(per ISO~21448 Clause~7) if rain exceeds 20~mm/h (reducing LiDAR point
density through absorption and scattering) or visibility falls below 200~m
(causing spurious reflections and reduced sensor range).

For each scene, we record: number of ground-truth objects, detections
(true positives, false positives, false negatives), mean geometric
divergence across ensemble members, and the number of uncertain detections
(divergence exceeding 0.5). These metrics map to the four evidence types
at G5 (Table~\ref{tab:evidence-types}).

% Requires: booktabs, amsmath, hyperref packages

\section{Methodology}\label{sec:methodology}

The analysis follows a five-step constructive integration methodology.
The input is a set of applicable standards; the output is an integrated
GSN argument pattern with identified inconsistencies and gaps.
Each step produces a verifiable intermediate artefact.

\subsection{System Definition and Standards Selection}\label{sec:system-def}

The case study component is a LiDAR-based 3D object detection module for
an SAE Level~4+ highly automated driving system. The architecture follows
the SECOND detector~\cite{Yan2018} with a deep ensemble of five
independently trained instances for uncertainty
quantification~\cite{Lakshminarayanan2017}.

The applicable standards are selected based on the following criteria:
(i)~the standard addresses at least one aspect of the component's safety
or cybersecurity assurance, (ii)~the standard is published (not draft),
and (iii)~the standard contains normative requirements or provides the
base assurance argument pattern. Five standards satisfy these criteria
(Table~\ref{tab:standards}). ISO/IEC~TR~5469 is included as informative
supplementary guidance.

The system receives an ASIL~D assignment from the hazard analysis and
risk assessment (HARA) per ISO~26262-3 Clause~6: severity S3 (life-threatening
injuries), exposure E4 (high probability in urban driving), controllability C3
(no human driver at SAE Level~4+). The threat analysis and risk assessment
(TARA) per ISO/SAE~21434 Clause~15 identifies three threat scenarios:
LiDAR spoofing (injected ghost points), adversarial point cloud perturbation,
and model poisoning through compromised training data.

\subsection{Step 1: Claim Extraction}\label{sec:step1}

For each applicable standard, we extract \emph{claims}: normative
requirements (\texttt{shall} statements), defined work products, and
explicitly required lifecycle activities that an assurance argument for
the AI-based perception component must address.

Each claim is assigned:
\begin{itemize}
    \item a unique identifier (e.g., \texttt{CLM-26262-HARA-01}),
    \item the source clause reference,
    \item a claim type (normative requirement, work product, or lifecycle activity),
    \item a target GSN goal node (determined in Step~3),
    \item a primary lifecycle phase (determined in Step~2).
\end{itemize}

Informative notes and recommendations are recorded but not treated as
normative claims. ISO/IEC~TR~5469, being a Technical Report, contributes
informative guidance only.

\subsection{Step 2: Lifecycle-Phase Mapping}\label{sec:step2}

Claims and clauses are mapped to six lifecycle phases that cover the
complete lifecycle of an AI-based perception component:
\begin{enumerate}
    \item \textbf{Concept / Requirements}: item definition, HARA, TARA,
          functional safety concept, SOTIF specification.
    \item \textbf{Design / Training}: AI model architecture, training process,
          cybersecurity control design, SOTIF design measures.
    \item \textbf{Verification \& Validation}: MC/DC coverage, scenario-based
          testing, AI-specific evaluation, penetration testing.
    \item \textbf{Integration / Deployment}: system integration, pre-deployment
          validation, release decision.
    \item \textbf{Operation / Monitoring}: runtime OOD detection, SOTIF field
          monitoring, cybersecurity incident monitoring.
    \item \textbf{Modification / Re-assurance}: OTA update management,
          change impact analysis, re-validation.
\end{enumerate}

The mapping produces a coverage matrix (Table~\ref{tab:coverage}) showing
how many clauses from each standard apply to each phase. Phases with low
coverage from any single standard are candidates for gaps.

\subsection{Step 3: GSN Construction}\label{sec:step3}

The integrated GSN argument pattern is constructed by:
\begin{enumerate}
    \item Starting with ISO/PAS~8800 Annex~B as the base structure. This
          provides goals G1--G6, strategy S1, and assumption A1.4.
    \item Adding context nodes: C1 (ASIL~D from HARA, ISO~26262-3 Clause~6)
          and C2 (TARA results, ISO/SAE~21434 Clause~15).
    \item Augmenting existing goals G2, G4, G5, G6 with claims from
          ISO~21448 and ISO/SAE~21434. The base pattern's G5 (V\&V)
          receives claims from all four normative standards.
    \item Adding new goals G7 (SOTIF residual risk acceptance, from
          ISO~21448 Clause~6.5) and G8 (cybersecurity risk management,
          from ISO/SAE~21434 Clause~3.1.11).
    \item Identifying G9 (modification assurance) as an undeveloped goal
          where no standard prescribes a complete re-assurance workflow
          for AI model modifications.
\end{enumerate}

The resulting structure (Table~\ref{tab:gsn-goals}) extends the base
pattern from 6 to 9 goals. Strategy S1 is reformulated to argue over
four assurance domains rather than one.

\subsection{Step 4: Junction-Point Analysis}\label{sec:step4}

A \emph{junction point} is a goal node where claims from two or more
standards meet. At each junction point, we analyse whether the standards
impose compatible requirements. Three types of inconsistency are
distinguished:
\begin{itemize}
    \item \textbf{Structural (S)}: different risk classification frameworks,
          evidence structures, or architectural assumptions.
    \item \textbf{Terminological (T)}: different terms for the same concept
          or the same term with different definitions.
    \item \textbf{Methodological (M)}: different methods prescribed for the
          same assurance objective.
\end{itemize}

The analysis identifies seven inconsistencies (Table~\ref{tab:inconsistencies}).
The central finding is I-2 (evidence type asymmetry at G5): four fundamentally
different evidence types---MC/DC structural coverage (ISO~26262), scenario-based
testing (ISO~21448), AI uncertainty quantification (ISO/PAS~8800), and
penetration testing (ISO/SAE~21434)---must support a single sufficiency
claim, yet no standard defines how they combine.

\subsection{Step 5: Gap Identification}\label{sec:step5}

Gaps are identified where the integrated argument requires a claim or
evidence that no applicable standard provides. Three gap types are
distinguished:
\begin{itemize}
    \item \textbf{Missing Claim (MC)}: the argument requires a claim
          but no standard contains it.
    \item \textbf{Missing Evidence (ME)}: a claim exists but no standard
          prescribes how to produce evidence for AI components.
    \item \textbf{Unresolved Inconsistency (UI)}: an inconsistency from
          Step~4 requires resolution that no standard provides.
\end{itemize}

Gaps are further classified as \emph{integration-induced} if they appear
only when standards are combined (not visible when each standard is applied
independently). The analysis identifies six gaps (Table~\ref{tab:gaps}),
of which two (Gap-3 and Gap-4) are integration-induced.

\subsection{Evaluation Protocol}\label{sec:evaluation}

The case study evaluation uses the CARLA simulator~\cite{Dosovitskiy2017}
under parametrically controlled weather conditions. Rain intensity varies
across five levels (0, 10, 25, 50, 100~mm/h) and fog visibility across
five levels (500, 200, 100, 50, 10~m), producing a grid of 25 weather
combinations. Each combination is evaluated over 50 scenes with randomised
traffic configurations (seed: 42).

A weather condition is classified as a \emph{SOTIF triggering condition}
(per ISO~21448 Clause~7) if rain exceeds 20~mm/h (reducing LiDAR point
density through absorption and scattering) or visibility falls below 200~m
(causing spurious reflections and reduced sensor range).

For each scene, we record: number of ground-truth objects, detections
(true positives, false positives, false negatives), mean geometric
divergence across ensemble members, and the number of uncertain detections
(divergence exceeding 0.5). These metrics map to the four evidence types
at G5 (Table~\ref{tab:evidence-types}).
